\documentclass[main.tex]{subfiles}
\begin{document}
\section{Conclusioni}\label{sec:cap 5}

L’attacco buffer overflow realizzato in questo elaborato ha permesso di comprendere nel dettaglio il funzionamento di una delle vulnerabilità più storiche e pericolose nel campo della sicurezza informatica. Attraverso lo studio del comportamento dello stack, l’analisi degli indirizzi di memoria, e la costruzione manuale di un payload con shellcode, è stato possibile sfruttare con successo una falla presente in un programma vulnerabile.
Questa attività ha evidenziato quanto sia cruciale l’adozione di meccanismi di protezione come l’ASLR, la non-eseguibilità dello stack e il canary stack protector, oggi presenti nei moderni sistemi operativi.





\end{document}